% problem_id: S001-lemma-relative-map-approximation
% source_id: S001 (Stacks Project)
% source_file: more-morphisms.tex
% statement_locator: more-morphisms.tex lines 11288--11364
% proof_locator: more-morphisms.tex lines 11366--11527
% env: lemma
% id_note: label 在本来源内唯一
% pairing: tight (verified)
% status: complete (statement + author proof verbatim + reason + locators)
%
\begin{lemma}
\label{lemma-relative-map-approximation}
Consider a diagram
$$
\vcenter{
\xymatrix{
X \ar[d] & Y \ar[d] \\
S & T \ar[l]
}
}
\quad\text{with points}\quad
\vcenter{
\xymatrix{
x \ar[d] & y \ar[d] \\
s & t \ar[l]
}
}
$$
where $S$ be a locally Noetherian scheme and the morphisms are
locally of finite type. Assume $\mathcal{O}_{S, s}$ is a G-ring.
Assume further we are given a local $\mathcal{O}_{S, s}$-algebra map
$$
\sigma : \mathcal{O}_{T, t} \longrightarrow \mathcal{O}_{S, s}^\wedge
$$
and a local $\mathcal{O}_{S, s}$-algebra map
$$
\varphi :
\mathcal{O}_{X, x}
\longrightarrow
\mathcal{O}_{Y_\sigma, y_\sigma}^\wedge
$$
where $Y_\sigma = Y \times_{T, \sigma} \Spec(\mathcal{O}_{S, s}^\wedge)$
and $y_\sigma$ is the unique point of $Y_\sigma$ lying over $y$.
For every $N \geq 1$ there exists a commutative diagram
$$
\xymatrix{
X \ar[d] & X \times_S V \ar[l] \ar[rd] &
W \ar[l]^-f \ar[r] \ar[d] &
Y \times_{T, \tau} V \ar[r] \ar[ld] & Y \ar[d] \\
S & & V \ar[ll] \ar[rr]^\tau & & T
}
$$
of schemes over $S$ and points $w \in W$, $v \in V$ such that
\begin{enumerate}
\item $v \mapsto s$, $\tau(v) = t$, $f(w) = (x, v)$, and $w \mapsto (y, v)$,
\item $(V, v) \to (S, s)$ is an elementary \'etale neighbourhood,
\item the diagram
$$
\xymatrix{
\mathcal{O}_{S, s}^\wedge \ar[r] & \mathcal{O}_{V, v}^\wedge \\
\mathcal{O}_{T, t} \ar[r]^{\tau^\sharp_v} \ar[u]_\sigma &
\mathcal{O}_{V, v} \ar[u]
}
$$
commutes module $\mathfrak m_v^N$,
\item $(W, w) \to (Y \times_{T, \tau} V, (y, v))$ is an
elementary \'etale neighbourhood,
\item the diagram
$$
\xymatrix{
\mathcal{O}_{X, x} \ar[r]_\varphi &
\mathcal{O}_{Y_\sigma, y_\sigma}^\wedge \ar[r] &
\mathcal{O}_{Y_\sigma, y_\sigma}/\mathfrak m_{y_\sigma}^N \ar@{=}[r] &
\mathcal{O}_{Y \times_{T, \tau} V, (y, v)}/\mathfrak m_{(y, v)}^N
\ar[d]_{\cong} \\
\mathcal{O}_{X, x} \ar[r] \ar@{=}[u] &
\mathcal{O}_{X \times_S V, (x, v)} \ar[r]^{f^\sharp_w} &
\mathcal{O}_{W, w} \ar[r] &
\mathcal{O}_{W, w}/\mathfrak m_w^N
}
$$
commutes. The equality comes from the fact that
$Y_\sigma$ and $Y \times_{T, \tau} V$ are canonically isomorphic over
$\mathcal{O}_{V, v}/\mathfrak m_v^N = \mathcal{O}_{S, s}/\mathfrak m_s^N$
by parts (2) and (3).
\end{enumerate}
\end{lemma}

\begin{proof}
After replacing $X$, $S$, $T$, $Y$ by affine open subschemes we
may assume the diagram in the statement of the lemma comes
from applying $\Spec$ to a diagram
$$
\vcenter{
\xymatrix{
A & B \\
R \ar[u] \ar[r] & C \ar[u]
}
}
\quad\text{with primes}\quad
\vcenter{
\xymatrix{
\mathfrak p_A & \mathfrak p_B \\
\mathfrak p_R \ar@{-}[u] \ar@{-}[r] & \mathfrak p_C \ar@{-}[u]
}
}
$$
of Noetherian rings and finite type ring maps.
In this proof every ring $E$ will be a Noetherian $R$-algebra endowed with
a prime ideal $\mathfrak p_E$ lying over $\mathfrak p_R$ and all ring maps
will be $R$-algebra maps compatible with the given primes. Moreover,
if we write $E^\wedge$ we mean the completion of the localization
of $E$ at $\mathfrak p_E$. We will also use without further mention
that an \'etale ring map $E_1 \to E_2$ such that
$\kappa(\mathfrak p_{E_1}) = \kappa(\mathfrak p_{E_2})$ induces
an isomorphism $E_1^\wedge = E_2^\wedge$ by
More on Algebra, Lemma \ref{more-algebra-lemma-flat-unramified}.

\medskip\noindent
With this notation $\sigma$ and $\varphi$ correspond to ring maps
$$
\sigma : C \to R^\wedge
\quad\text{and}\quad
\varphi : A \longrightarrow (B \otimes_{C, \sigma} R^\wedge)^\wedge
$$
Here is a picture
$$
\xymatrix{
A \ar@/^1em/[rrr]^\varphi &
B \ar[r] &
B \otimes_{C, \sigma} R^\wedge \ar[r] &
(B \otimes_{C, \sigma} R^\wedge)^\wedge \\
R \ar[r] \ar[u] &
C \ar[r]^\sigma \ar[u] & R^\wedge \ar[u] \ar[ru]
}
$$
Observe that $R^\wedge$ is a G-ring by
More on Algebra, Proposition
\ref{more-algebra-proposition-Noetherian-complete-G-ring}.
Thus $B \otimes_{C, \sigma} R^\wedge$ is a G-ring by
More on Algebra, Proposition
\ref{more-algebra-proposition-finite-type-over-G-ring}.
By Lemma \ref{lemma-map-approximation} (translated into algebra)
there exists an \'etale ring map $B \otimes_{C, \sigma} R^\wedge \to B'$
inducing an isomorphism
$\kappa(\mathfrak p_{B \otimes_{C, \sigma} R^\wedge})
\to \kappa(\mathfrak p_{B'})$
and an $R$-algebra map $A \to B'$ such that the composition
$$
A \to B' \to (B')^\wedge = (B \otimes_{C, \sigma} R^\wedge)^\wedge
$$
is the same as $\varphi$ modulo
$(\mathfrak p_{(B \otimes_{C, \sigma} R^\wedge)^\wedge})^N$.
Thus we may replace $\varphi$ by this composition because the
only way $\varphi$ enters the conclusion is via the commutativity
requirement in part (5) of the statement of the lemma.
Picture:
$$
\xymatrix{
& & B' \ar[r] & (B')^\wedge \ar@{=}[d] \\
A \ar[rru] &
B \ar[r] &
B \otimes_{C, \sigma} R^\wedge \ar[r] \ar[u] &
(B \otimes_{C, \sigma} R^\wedge)^\wedge \\
R \ar[r] \ar[u] &
C \ar[r]^\sigma \ar[u] & R^\wedge \ar[u] \ar[ru]
}
$$
Next, we use that $R^\wedge$ is a filtered colimit of smooth
$R$-algebras (Smoothing Ring Maps, Theorem \ref{smoothing-theorem-popescu})
because $R_{\mathfrak p_R}$ is a G-ring by assumption. Since $C$ is of
finite presentation over $R$ we get a factorization
$$
C \to R' \to R^\wedge
$$
for some $R \to R'$ smooth, see
Algebra, Lemma \ref{algebra-lemma-characterize-finite-presentation}.
After increasing $R'$ we may assume there exists
an \'etale $B \otimes_C R'$-algebra $B''$ whose base change
to $B \otimes_{C, \sigma} R^\wedge$ is $B'$, see
Algebra, Lemma \ref{algebra-lemma-etale}.
Then $B'$ is the filtered colimit of these $B''$ and we
conclude that after increasing $R'$ we may assume there is
an $R$-algebra map $A \to B''$ such that $A \to B'' \to B'$ is
the previously constructed map (same reference as above). Picture
$$
\xymatrix{
& & B'' \ar[r] & B' \ar[r] & (B')^\wedge \ar@{=}[d] \\
A \ar[rru] &
B \ar[r] &
B \otimes_C R' \ar[r] \ar[u] &
B \otimes_{C, \sigma} R^\wedge \ar[r] \ar[u] &
(B \otimes_{C, \sigma} R^\wedge)^\wedge \\
R \ar[r] \ar[u] &
C \ar[r] \ar[u] &
R' \ar[r] \ar[u] &
R^\wedge \ar[u] \ar[ru]
}
$$
and
$$
B' = B'' \otimes_{(B \otimes_C R')} (B \otimes_{C, \sigma} R^\wedge)
$$
This means that we may replace $C$ by $R'$, $\sigma : C \to R^\wedge$ by
$R' \to R^\wedge$, and $B$ by $B''$ so that we simplify to the diagram
$$
\xymatrix{
A \ar[r] &
B \ar[r] &
B \otimes_{C, \sigma} R^\wedge \\
R \ar[r] \ar[u] &
C \ar[r]^\sigma \ar[u] & R^\wedge \ar[u]
}
$$
with $\varphi$ equal to the composition of the horizontal arrows
followed by the canonical map from $B \otimes_{C, \sigma} R^\wedge$
to its completion.
The final step in the proof is to apply Lemma \ref{lemma-map-approximation}
(or its proof)
one more time to $\Spec(C)$ and $\Spec(R)$ over $\Spec(R)$ and the map
$C \to R^\wedge$. The lemma produces a ring map $C \to D$
such that $R \to D$ is \'etale, such that
$\kappa(\mathfrak p_R) = \kappa(\mathfrak p_D)$, and such that
$$
C \to D \to D^\wedge = R^\wedge
$$
is equal to $\sigma : C \to R^\wedge$ modulo $(\mathfrak p_{R^\wedge})^N$.
Then we can take
$$
V = \Spec(D)
\quad\text{and}\quad
W = \Spec(B \otimes_C D)
$$
as our solution to the problem posed by the lemma. Namely the diagram
$$
\xymatrix{
A \ar[r] &
B \otimes_{C, \sigma} R^\wedge \ar[r] &
B \otimes_{C, \sigma} R^\wedge/(\mathfrak p_{R^\wedge})^N \ar@{=}[r] &
B \otimes_C D/(\mathfrak p_D)^N \\
A \ar@{=}[u] \ar[r] &
A \otimes_R D \ar[r] &
B \otimes_R D \ar[r] &
B \otimes_C D/(\mathfrak p_D)^N \ar@{=}[u]
}
$$
commutes because $C \to D \to D^\wedge = R^\wedge$ is equal to
$\sigma$ modulo $(\mathfrak p_{R^\wedge})^N$. This proves part (5)
and the other properties are immediate from the construction.
\end{proof}
